\section{De los menores al logaritmo exacto}
En la máscara fija sea $\mathcal A_\rho$ el álgebra $z^U z^V=z^{U\cup V}$
de norma $\|\sum a_Uz^U\|_\rho=\sum |a_U|\rho^{|U|}$, $\rho\ge2$.
En una evaluación booleana se activan sólo los conflictos internos a los sitios activos.
Si $a_U=\sum_{Y\subset U}(-1)^{|U|-|Y|}{\bf1}_{Y\ {\rm independiente}}$
(suma inclusiva), la partición exacta es $Z=\sum_U a_Uz^U p_U$ por inversión.
Los ceros RC se omiten. $a_U$ factoriza por componentes diagonales y anula los
vértices aislados. Sea $w_U=a_Uz^U\prod b_e$, $\eta_j=\rho^2(\epsilon_A+1/j)$,
$S=j\eta_j$, con $\epsilon_A$ la máxima masa diagonal de $P$.
F01 compara $D$ y $P$ por balanceo, y los $b_e$ son comparables a $D_e$.

Por F04 un polímero anclado en $i$ de talla $v\ge2$ cuesta, tras sumar en la norma,
\[
 b_i(C_B\rho)^v\,\frac{v^{v-2}}{(v-1)!}(C\epsilon_A)^{v-1}.
\]
La constante incluye una reserva $e^{Bv}$. La cota sale al elegir un árbol
etiquetado de vecinos diagonales y sumar desde sus hojas. La masa por vecino es
$O(\epsilon_A)$, la masa por fibra RC es $O(1)$, y $\sum b_i=O(j)$.
Después de la suma en $v$ la cota es $C_B b_i\eta_j$; un polímero de talla $v$
tiene masa incompatible $O_B(v\eta_j)$ y la masa total es $O_B(S)$.
El criterio de Kotecký--Preiss con penalización $v$ da una serie de clústeres
absolutamente convergente para el logaritmo del gas, incluida la reserva.
En particular $\sum_U\|w_U\|_\rho C_0^{|U|}e^{B'|U|}\le e^{C_BS}$
para la reserva fija requerida, por suma de colecciones con actividades mayoradas.

Para $\xi:E\to\mathcal A_\rho$ lineal en un tubo fijo ponemos
$Z_\xi=\sum_U w_Ue^{\xi[T_U]}$, $G(\xi)=\log Z_\xi$ por la serie de clústeres,
$E_\xi f=Z_\xi^{-1}\sum_Uw_Ue^{\xi[T_U]}f(T_U)$ y $m=E_\xi T$.
Una y dos marcas cuestan $3v$ y $(3v)^2$ por clúster de talla total $v$.
Así $\|G\|,\|m\|_\pi,\|Dm\|\le CS$ y $\|Z_\xi^{\pm1}\|\le e^{CS}$.
Para evaluar $L_s$ en $E\widehat\otimes_\pi\mathcal A_\rho$ se polarizan sus
coeficientes en $E$ y se multiplican sus valores en el álgebra. En grado $l$ cuesta
a lo sumo $e^l$, absorbido al reducir el radio. En la bola reducida el Hessiano
sigue acotado por $C/j$. El mapa $\xi\mapsto DL_s(m(\xi))$ preserva una bola
de radio $C'(1+S)/j$ y tiene constante $CS/j<1/2$. Su punto fijo es real
en cada evaluación booleana, por iteración desde cero. F05 lo denota por $h$.

Los tensores centrados $M_l=E_\xi[(T-m)^{\otimes l}]$ se mayoran en el álgebra
simétrica con norma proyectiva. Marcar cada clúster antes de sumar da, para $0<w\le w_0$,
una norma del logaritmo de la generatriz centrada de a lo sumo $CSw^2$:
la parte a partir de dos marcas cuesta $Cw^2v^2e^{3w_0v}$ por clúster.
Así $\|M_l\|\le l!w^{-l}e^{CSw^2}\le[C(\sqrt{lS}+l)]^l$, eligiendo
$w=\min(w_0,\sqrt{l/(1+S)})$. En $B(Y)=\exp(L_s(m+Y)-L_s(m)-\xi[Y])$
el Hessiano da $\|B(Y)\|\le\exp(C\|Y\|_\pi^2/j)$ en una bola de radio de orden $j$.
El Taylor $b_l$ tiene $b_0=1,b_1=0$ y
$\|b_l\|\le(C/\sqrt{jl})^l$ hasta $H=\lfloor\kappa j\rfloor$ por Cauchy.
La contracción es legítima: las formas extendidas son multilineales sobre el álgebra,
y sus restricciones a $E^l$ se contraen con los tensores de momentos multiplicando
el coeficiente del álgebra. La norma proyectiva evita factores de dimensión.
Al separar dos factores se obtiene, para $q=C(\sqrt{\eta_j}+\sqrt\kappa)<1$,
\[
 \sum_{2\le l\le H}\|b_l:M_l\|\le
 \frac Cj\sum_{l\ge2}(S+l)q^{l-2}\le C'(1+S)/j.
\]

Sea $P_H$ el Taylor truncado y $D_\xi=Z_\xi\exp(L_s(m)-\xi[m])$.
Las normas de $D_\xi^{\pm1}$ cuestan $e^{CS}$. Para $|U|\le aj$ se reemplaza
$p_U$ con error relativo $Ce^{-cj}$ y se factoriza $e^{L_s(T_U)}$
como $e^{L_s(m)-\xi[m]} e^{\xi[T_U]} B(T_U-m)$.
En radio $R=c_1\sqrt{\kappa}\,j$, $\|B\|\le e^{CR^2/j}$.
Elegimos $a\ll\sqrt\kappa$ y $\eta_F$ aún menor para que el resto de Cauchy
en $\|T_U-m\|\le3aj+CS$ sea $O(e^{-c_2\kappa j})$.
Para $|U|>aj$, el mayorante global $C_0^{|U|}$ y una reserva $e^{|U|}$
dan $e^{-aj+CS}$, incluso tras normalizar. Para la cola artificial del polinomio,
$l!\|b_l\|\le(C\sqrt\kappa)^l$, de donde
$\|P_H(T_U-m)\|\le\exp(C\sqrt\kappa(3|U|+CS))$, dentro de la misma reserva.
El orden es: reservar $a_*,C_0$ y un $B$ fijo que pague el campo, $C_0$ y las
exponenciales por sitio; reducir $\kappa$, después $a$, después $\eta_F$.
Se puede acotar $\sqrt\kappa\le1$ al fijar $B$. El coste $e^{CS}$ es absorbido
por cada cola al elegir $\eta_F$ tras los exponentes positivos anteriores.
Así $\|D_\xi^{-1}Z-E_\xi P_H(T-m)\|\le Ce^{-c'j}$. Resulta para $j$ grande
\[
 \|D_\xi^{-1}Z-1\|_\rho\le C(1+S)/j+Ce^{-c'j}<1/2.
\]
En cada evaluación $Z_\xi=\exp G(\xi)>0$ por realidad, luego $D_\xi>0$ y $Z>0$.
En la activación total $Z=Q(B)/\operatorname{per}A$, con $\operatorname{per}A>0$.
Sólo ahora se toma el logaritmo físico. La serie logarítmica y los tres factores
dan $\|F\|_\rho\le C(1+j\rho^2(\epsilon_A+1/j))$. La deducción es determinista,
sirve en $H_j$ y fija $\eta_F$ mediante los márgenes antes de conocer $\rho$ y $b$.

El umbral físico de prescripciones cortas cumple $a\le a_*$. La constante del truncamiento
$H=\lfloor\kappa_T j\rfloor$ se fija tras la reserva, y antes de $a$ y de $\eta_F$.
